%%%%%%%%% pl7.tex %%%%%%%%

\bs\bs\ni{\large\bf 3. General Structure Theory}

\medskip\ni{\bf 3.1. Algebra structure, Markov trace}

The proof of Theorem 3.1.3 below is routine for those conversant
with [J1] or [GHJ]. We include it since, as stated, it can
be useful in determining principal graphs. Recall that in a planar
algebra $P$, \ $e_k$ denotes the idempotent in $P_k$ equal to
$\frac{1}{\delta}(| |\dots \quad )$.

{\bf Lemma 3.1.1} Let $P$ be a finite-dimensional
spherical nondegenerate planar algebra over an algebraically closed
field. Then for each  $k$, \
$P_{k-1}e_{k-1}P_{k-1}$ is a 2-sided ideal, denoted $I_k$, in $P_k$
and if ${\cal M}_k$ is a set of minimal idempotents in $P_k$
generating all the distinct minimal ideals in $P_k/I_k$, we have
\begin{verse}
(i)   $pP_ke_{k-1}=0$ \ for $p$ in ${\cal M}_k$

(ii)  $pP_kq=0$ for $p\neq q$ in ${\cal M}_k$

(iii) For each $x$ in $P_k\backslash I_k$ there is a
$p\in\cal M_k$ with $xP_kp\neq 0$.

(iv) ${\op{tr}}(p)\neq 0$ for all $p\in {\cal M}_k$.

(v) $I_{k+2}=\bigoplus_{p\in{\cal M}_k}
P_{k+2}pe_{k+1}P_{k+2}$, \ $pe_{k+1}$ being a minimal idempotent in
$P_{k+2}$.
\end{verse}
Moreover, if, for each $k$, $\cal N_k$ is a set of minimal
idempotents of $P_k$ satisfying ${\op{(i)}}\dots{\op{(iv)}}$ (with
${\cal M}_k$ replaced by ${\cal N}_k$), then there is an invertible
$u_k$ in $P_k$ with $u_k\cal N_ku^{-1}_k={\cal M}_k$ (so in
particular (v) is true for ${\cal N}_k$).

\bigskip
{\sc Proof.} To see that $P_{k-1}e_{k-1}P_{k-1}$ is an
ideal, consider the maps $\a,\beta: P_k\to P_{k-1}$ given by the
annular tangles
\[
	\begin{picture}(0,0)%
\epsfig{file=ps/pic68.ps}%
\end{picture}%
\setlength{\unitlength}{0.00041700in}%
%
\begingroup\makeatletter\ifx\SetFigFont\undefined
% extract first six characters in \fmtname
\def\x#1#2#3#4#5#6#7\relax{\def\x{#1#2#3#4#5#6}}%
\expandafter\x\fmtname xxxxxx\relax \def\y{splain}%
\ifx\x\y   % LaTeX or SliTeX?
\gdef\SetFigFont#1#2#3{%
  \ifnum #1<17\tiny\else \ifnum #1<20\small\else
  \ifnum #1<24\normalsize\else \ifnum #1<29\large\else
  \ifnum #1<34\Large\else \ifnum #1<41\LARGE\else
     \huge\fi\fi\fi\fi\fi\fi
  \csname #3\endcsname}%
\else
\gdef\SetFigFont#1#2#3{\begingroup
  \count@#1\relax \ifnum 25<\count@\count@25\fi
  \def\x{\endgroup\@setsize\SetFigFont{#2pt}}%
  \expandafter\x
    \csname \romannumeral\the\count@ pt\expandafter\endcsname
    \csname @\romannumeral\the\count@ pt\endcsname
  \csname #3\endcsname}%
\fi
\fi\endgroup
\begin{picture}(7845,2175)(2079,-2483)
\put(5701,-1561){\makebox(0,0)[lb]{\smash{\SetFigFont{12}{14.4}{rm}and}}}
\end{picture}

\]
\ni respectively. A  diagram shows that
$xe_{k-1}y=\a(x)e_{k-1}\beta(y)$ for $x,y\in P_k$.

By Corollary 1.30, $P_k$ is semisimple and multimatrix since $K$ is
algebraically closed. Thus properties (i), (ii), (iii) and (v) are
obvious for ${\cal M}_k$. If $p\neq 0$ satisfied tr$(p)=0$, then tr
would vanish on the whole matrix algebra containing $p$ which
would then be orthogonal to $P_k$.

Finally, suppose we are given ${\cal N}_k$ satisfying (i)--(iv).
Then since $P_k$ is multimatrix, each $p$ in  ${\cal N}_k$ belongs to
a unique matrix algebra summand in which there is an invertible
$u_p$ with $u_ppu_p^{-1}\in{\cal M}_k$. Putting together the
$u_p$'s, and the identity of $I_k$, we get $u_k$. Property (v) for
${\cal N}_k$ then follows from (iii). \qed



\bs{\bf Definition 3.1.2.} With $P$ and ${\cal N}_k$ as in 3.11, we
define the {\it principal graph} $\G_P$ of $P$ to be the
(bipartite) graph whose vertices are $\bigcup_{k\geq 0}\cal N_k$
with distinguished vertex $*$ so that ${\cal N}_0=\{*\}$, and
dim $(pP_{k+1}q)$ edges between $p\in{\cal N}_k$ and $q\in{\cal N}_{k+1}$. 
Let $d_p$ denote the distance from $p$ to $*$ on $\G_P$.

\bs\ni
{\bf Theorem 3.1.3} As an algebra, $P_k$ is isomorphic to the
algebra whose basis is random walks of length $2k$ on $\G_P$
beginning and ending at $*$ with multiplication rule
$w_1w_2=w_3$ if the first half of the walk $w_2$ is equal to the
second half of $w_1$, and $w_3$ is the first half of $w_1$ followed
by the second half of $w_2$; \ 0 otherwise. Moreover, if $\vec t$
is the function from the vertices of $\G_P$ to $K$, \
$\vec t_p=\d^{d_p}{\op{tr}}(p)$, \ $\vec t$ is an eigenvector for
the adjacency matrix of $\G$, eigenvalue $\d$.

\bigskip
{\sc Proof.} The first assertion is easily equivalent to showing
that the Bratteli diagram (see [GHJ]) of the multimatrix
algebra $P_k$ in $P_{k+1}$ is the bipartite graph consisting of
those vertices $p$ with $d_p\equiv k$ (mod 2) and $d_p\leq k$,
connected to those with
$d_p\equiv (k+1)$ (mod 2) and $d_p\leq k+1$ with appropriate
multiplicities. Observe first that $P_0=\Bbb C$ and ${\cal M}_1$ is a
set of minimal projections, one for each matrix algebra summand of
$P_1$, so the Bratelli diagram is correct for $P_0\subset P_1$. Now
proceed by induction on $k$. The trace on $P_k$ is nondegenerate,
as is its restriction to $P_{k-1}$ so one may perform the abstract
``basic construction" of [J1] to obtain the algebra
$\<P_k,e_{P_{k-1}}\>$ which is multimatrix and isomorphic to
$P_k\otimes_{P_{k-1}} P_k$ as a $P_k-P_k$ bimodule via the map
$x\otimes y\mapsto xe_{P_{k-1}} y$. Moreover the matrix algebra
summands of $\<P_k,e_{P_{k-1}}\>$ are indexed by those of $P_{k-1}$,
which by induction are the vertices of $\G_P$ with $d_p\leq k-1$, \
$d_P\equiv (k+1)$ mod 2.  If one defines the trace tr on
$\<P_k,e_{P_{k-1}}\>$ by $\widetilde{\op{tr}}(xe_{P_{k-1}}y)=
\frac{1}{\d^2} {\op{tr}}(xy)$ then the traces of minimal
projections in $\<P_k,e_{P_{k-1}}\>$ are $\frac{1}{\d^2}$ times
those in $P_{k-1}$. Moreover, setting $\g(xe_{P_{k-1}}y)=xe_ky$
defines an algebra homomorphism from $\<P_k,e_{P_{k-1}}\>$ which is
injective by property (iv) and onto $I_{k+1}$. And
Tr$(xe_ky)=\frac 1\d$ Tr$(xy)$ so tr
$=\widetilde{\op{tr}}\circ\g^{-1}$ on $I_{k+1}$. Properties (i),
(ii) and (iii) ensure that the other  vertices of the Bratteli
diagram for $P_k\subset P_{k+1}$ are labelled by vertices $p$ of
$\G_P$ with $d_P=k+1$. And the number of edges on $\G_P$ connecting
a $p$ in ${\cal M}_k$ to a $q$ in ${\cal M}_{k+1}$ is by definition the
number of edges in the Bratteli diagram. That there are no edges
between ${\cal M}_{k+1}$ and ${\cal M}_j$, $j < k$, follows from
$$
(xe_{k-1}y)p=\frac{1}{\d^2} \ x(e_{k-1}e_ke_{k-1})yp
= \frac{1}{\d^2} \ x(e_{k-1}e_kpe_{k-1})y=0
$$
by (i) for $x,y\in P_{k-1}$ and $p\in {\cal M}_{k+1}$.

Finally, the (normalized) trace of a minimal projection $p$ in
$I_k$ is $\d^{d_P-k}\vec t_p$ so the assertion about the trace
follows as usual (see [J1]). \qed

\bs{\bf Remarks.} (1) Similarly, the algebras $P_{1,k}$ have a
principal graph $\G'_P$ with the trace vector $\vec s$. We call
$\G'_P$ the dual principal graph. Ocneanu has shown, in the $C^*$
case,  how to associate numerical data encoding the ensuing
embedding of the random walk algebra of $\G'_P$ into that of
$\G_P$. This completely captures the planar algebra structure and
is analogous to choosing local coordinates on a manifold. The same
could be done under the hypotheses of Lemma 3.1.1. The principal
graphs alone do not determine the planar algebra --- for instance
the algebras $(P^\s)^{\z/4\z}$ and
$(P^\s)^{\z/2\z\oplus\z/2\z}$ of 2.8 have the same principal graph
but are readily distinguished by counting fixed points under the
rotation.

(2) In fact, the assumption of nondegeneracy on $P$ in 3.1.1 and
3.1.3 could be replaced by the hypothesis $P_k/I_k$ semisimple.
Then conditions (i)--(iv) could be  used to inductively guarantee
nondegeneracy.

(3) If $P$ had been a $C^*$-planar algebra, we would have a
theorem (3.1.3)${}^*$ with all tr$(p)$ positive, all $p$'s
projections, and the obvious adjoint \ ${}^*$ \  on random walks.

\bs
Theorem 3.1.3 can be used to compute the principal graphs for
Temperley Lieb and the Fuss Catalan algebra (it is the ``middle
pattern" method of [BJ2]). We now illustrate its use by
calculating the principal graph of the nondegenerate planar
algebras coming from Example 2.2. We work over $\Bbb C$ for
convenience, and in the
$C^*$ case to simplify life.


Let $(A,{\op{TR}})$ be a finite-dimensional unital $C^*$-algebra
with normalized faithful (positive) trace TR. The labelling set $L$
is $L_1=A$. We choose a number $\d > 0$ and let $\tau_p$ be TR$(p)$
for projections $p\in A$. A labelled network is then a disjoint
union of smoothly embedded circles, each one containing a (possibly
empty) sequence of 1-boxes labelled by elements in $A$. We define
the partition function $Z$ of such a collection of circles to be
$\d^{\# \ {\text{(circles)}}}\prod_{\text{(circles)}}
{\op{TR}}(a_1a_2\dots a_n)$, where
$a_1a_2\dots a_n$ are the labels  on the given circle, numbered in
order around the circle. The partition function $Z$ is obviously
multiplicative so we define $P^{(A,{\op{TR}})}$ to be the
nondegenerate planar algebra, with obvious $*$-structure, defined
by 1.23. It is linearly spanned by Temperley-Lieb diagrams with a
single labelled box on each string. The relations of Example 2.2
hold, noting that
$Z(\blrighto{p})=\d\tau_p$. For certain values of $\d$ and traces
TR we will compute the principal (and dual principal) graphs of
$P^{A,{\op{TR}})}$ and the Markov trace, and show it to be a planar
$C^*$-algebra. Let us first describe the graphs. Let ${\cal M}=\{p\}$
be a set of minimal projections in $A$, one for each matrix algebra
direct summand and let
$n_p={\op{dim}}(pA)$. Let $\frak S({\cal M})$ be the free
semigroup with identity on ${\cal M}$. Let $\l:{\op{Proj}}\to\Bbb
N\cup\{\infty\}$ be a function, and $W_\l$ be the set of words in
 $\frak S({\cal M})$  which contain no consecutive string of $p$'s
longer than $\l(p)$, for each $p$.

\medskip{\bf Definition 3.1.4.} The graph $\G_{A,\l}$ is the rooted
tree having vertices $W_\l$, with $n_p$ edges between $w$ and
$wp$ for every $p\in\cal M$ with $\{w\cup wp\}\subset W_f$. The
root $*$ is the identity of $\frak S({\cal M})$.

\medskip Thus if $\l(p)=1$ for all $p$ and $A$ is abelian,
$\G_{A,\l}$ is the regular tree of valence $|{\cal M}|$. If
$\l(p)=\infty$ for all
$p$ and $A$ is abelian, the root $*$ of the tree $\G_{A,\l}$ has
valence
$|{\cal M}|$ and all other vertices have valence $|{\cal M}|+1$.
If $A=\Bbb Cp+\Bbb Cq$ and $\l(p)=1$, \ $\l(q)=2$, the tree
$\G_{A,\l}$ is as in Figure~\ref{pic69}.
%%%%%%%%%%%%%%%%%%%%%%%%%%%%%%%%%%%%%%%%%%%%%%%%%%%%%%%%%%%%%%%%
\[
	\begin{picture}(0,0)%
\epsfig{file=ps/pic69.ps}%
\end{picture}%
\setlength{\unitlength}{0.00041700in}%
%
\begingroup\makeatletter\ifx\SetFigFont\undefined
% extract first six characters in \fmtname
\def\x#1#2#3#4#5#6#7\relax{\def\x{#1#2#3#4#5#6}}%
\expandafter\x\fmtname xxxxxx\relax \def\y{splain}%
\ifx\x\y   % LaTeX or SliTeX?
\gdef\SetFigFont#1#2#3{%
  \ifnum #1<17\tiny\else \ifnum #1<20\small\else
  \ifnum #1<24\normalsize\else \ifnum #1<29\large\else
  \ifnum #1<34\Large\else \ifnum #1<41\LARGE\else
     \huge\fi\fi\fi\fi\fi\fi
  \csname #3\endcsname}%
\else
\gdef\SetFigFont#1#2#3{\begingroup
  \count@#1\relax \ifnum 25<\count@\count@25\fi
  \def\x{\endgroup\@setsize\SetFigFont{#2pt}}%
  \expandafter\x
    \csname \romannumeral\the\count@ pt\expandafter\endcsname
    \csname @\romannumeral\the\count@ pt\endcsname
  \csname #3\endcsname}%
\fi
\fi\endgroup
\begin{picture}(7244,4278)(1479,-4207)
\put(4471,-181){\makebox(0,0)[lb]{\smash{\SetFigFont{12}{14.4}{rm}$*$}}}
\end{picture}

\]
\begin{center}
	Figure 3.1.5
\end{center}
%%%%%%%%%%%%%%%%%%%%%%%%%%%%%%%%%%%%%%%%%%%%%%%%%%%%%%%%%%%
Recall the polynomials $T_n(x)$ of [J1], $T_1\!=\!1$,
$T_2\!=\!1$,
$T_{n+1}=T_n-xT_{n-1}$ and the ``Jones-Wenzl" projections
$f_k\in TL(k)$ with $f^*_k=f^2_k=f_k$, \
$f_ke_i\!=\!0$ for $i\!=\!1,2,\dots ,k-1$ so that $f_k$ are elements
of any (spherical) planar algebra, and
tr$(f_k)=T_{k-2}(\frac{1}{\d^2})$ if $T_j(\frac{1}{\d^2})\neq 0$
for $j < k+2$.

{\bf Theorem 3.1.6} With notation as above, suppose
$\tau_p\d=2\cos \pi/(\l(p)+2)$ ($\tau_p\d\geq 2$ if $\l(p)=\infty$).
Then $P^{(A,{\op{TR}})}$ is a (spherical) $C^*$-planar algebra with
principal and dual principal graphs equal to $\G_{A,\l}$. The
(normalized) trace of the minimal projection in
$P_k^{(A,{\op{TR}})}$ corresponding to the word
$w\!=\!p_1^{m_1}p_2^{m_2}\dots p_r^{m_r}$ (\,$\sum m_i\!=\!k$ and
$p_i\neq p_{i+1}$) is
$\prod^r_{i=1} \tau_{p_i}T_{m_i}(\tau_{p_i}\d)$.

\bigskip
{\sc Proof.} We shall give explicit projections satisfying
conditions (i)--(iv) of 3.1.1. The key observation is that tangles
with a fixed $p$ labelling each string form a Temperley-Lieb
subalgebra $B_p$ with parameter $\d\tau_{p}$ (and identity
$\boxf{p} \ \boxf{p}\dots \boxf{p}$\,). So if $m <\l(p)$ we
consider the projection
\[
	\begin{picture}(0,0)%
\epsfig{file=ps/pic70.ps}%
\end{picture}%
\setlength{\unitlength}{0.00041700in}%
%
\begingroup\makeatletter\ifx\SetFigFont\undefined
% extract first six characters in \fmtname
\def\x#1#2#3#4#5#6#7\relax{\def\x{#1#2#3#4#5#6}}%
\expandafter\x\fmtname xxxxxx\relax \def\y{splain}%
\ifx\x\y   % LaTeX or SliTeX?
\gdef\SetFigFont#1#2#3{%
  \ifnum #1<17\tiny\else \ifnum #1<20\small\else
  \ifnum #1<24\normalsize\else \ifnum #1<29\large\else
  \ifnum #1<34\Large\else \ifnum #1<41\LARGE\else
     \huge\fi\fi\fi\fi\fi\fi
  \csname #3\endcsname}%
\else
\gdef\SetFigFont#1#2#3{\begingroup
  \count@#1\relax \ifnum 25<\count@\count@25\fi
  \def\x{\endgroup\@setsize\SetFigFont{#2pt}}%
  \expandafter\x
    \csname \romannumeral\the\count@ pt\expandafter\endcsname
    \csname @\romannumeral\the\count@ pt\endcsname
  \csname #3\endcsname}%
\fi
\fi\endgroup
\begin{picture}(3805,4258)(2379,-4297)
\put(4126,-2236){\makebox(0,0)[lb]{\smash{\SetFigFont{12}{14.4}{rm}$f_m$}}}
\put(3676,-3736){\makebox(0,0)[lb]{\smash{\SetFigFont{12}{14.4}{rm}$p$}}}
\put(2626,-3736){\makebox(0,0)[lb]{\smash{\SetFigFont{12}{14.4}{rm}$p$}}}
\put(5776,-3736){\makebox(0,0)[lb]{\smash{\SetFigFont{12}{14.4}{rm}$p$}}}
\put(3676,-736){\makebox(0,0)[lb]{\smash{\SetFigFont{12}{14.4}{rm}$p$}}}
\put(4726,-736){\makebox(0,0)[lb]{\smash{\SetFigFont{12}{14.4}{rm}$p$}}}
\put(5776,-736){\makebox(0,0)[lb]{\smash{\SetFigFont{12}{14.4}{rm}$p$}}}
\put(2626,-736){\makebox(0,0)[lb]{\smash{\SetFigFont{12}{14.4}{rm}$p$}}}
\end{picture}

\]
\ni where $f_m$ is calculated in Temperley-Lieb with $j$ strings
and $Z(\caright\ )=\tau_p\d$. Now set
\[
	\begin{picture}(0,0)%
\epsfig{file=ps/pic71.ps}%
\end{picture}%
\setlength{\unitlength}{0.00041700in}%
%
\begingroup\makeatletter\ifx\SetFigFont\undefined
% extract first six characters in \fmtname
\def\x#1#2#3#4#5#6#7\relax{\def\x{#1#2#3#4#5#6}}%
\expandafter\x\fmtname xxxxxx\relax \def\y{splain}%
\ifx\x\y   % LaTeX or SliTeX?
\gdef\SetFigFont#1#2#3{%
  \ifnum #1<17\tiny\else \ifnum #1<20\small\else
  \ifnum #1<24\normalsize\else \ifnum #1<29\large\else
  \ifnum #1<34\Large\else \ifnum #1<41\LARGE\else
     \huge\fi\fi\fi\fi\fi\fi
  \csname #3\endcsname}%
\else
\gdef\SetFigFont#1#2#3{\begingroup
  \count@#1\relax \ifnum 25<\count@\count@25\fi
  \def\x{\endgroup\@setsize\SetFigFont{#2pt}}%
  \expandafter\x
    \csname \romannumeral\the\count@ pt\expandafter\endcsname
    \csname @\romannumeral\the\count@ pt\endcsname
  \csname #3\endcsname}%
\fi
\fi\endgroup
\begin{picture}(4095,3345)(2229,-3384)
\put(2551,-736){\makebox(0,0)[lb]{\smash{\SetFigFont{12}{14.4}{rm}$p_1$}}}
\put(3901,-736){\makebox(0,0)[lb]{\smash{\SetFigFont{12}{14.4}{rm}$p_2$}}}
\put(5701,-736){\makebox(0,0)[lb]{\smash{\SetFigFont{12}{14.4}{rm}$p_r$}}}
\put(2551,-2911){\makebox(0,0)[lb]{\smash{\SetFigFont{12}{14.4}{rm}$p_1$}}}
\put(3901,-2911){\makebox(0,0)[lb]{\smash{\SetFigFont{12}{14.4}{rm}$p_2$}}}
\put(5701,-2911){\makebox(0,0)[lb]{\smash{\SetFigFont{12}{14.4}{rm}$p_r$}}}
\put(2401,-1786){\makebox(0,0)[lb]{\smash{\SetFigFont{12}{14.4}{rm}$f_{m_1}$}}}
\put(3751,-1786){\makebox(0,0)[lb]{\smash{\SetFigFont{12}{14.4}{rm}$f_{m_2}$}}}
\put(5551,-1861){\makebox(0,0)[lb]{\smash{\SetFigFont{12}{14.4}{rm}$f_{m_r}$}}}
\end{picture}

\]
\ni where we have combined the $m_j$ strings at the top (and
bottom) of $f_{m_j}$ into one. Orientations are completely
forgotten and may be inserted, if required, so as to satisfy
Definition 1.7.

Condition (1) of 3.1.1 is easy: $P_k$ is linearly spanned by
Temperley Lieb diagrams with matrix units  (with the $p$'s among
the diagonal ones) from the simple summands of $A$ in a single box
on each string. If $x$ is such an element, then if
$xe_{k-1}$ is non-zero, the product $p_wxe_{k-1}$ contains
\[
	\begin{picture}(0,0)%
\epsfig{file=ps/pic72.ps}%
\end{picture}%
\setlength{\unitlength}{0.00041700in}%
%
\begingroup\makeatletter\ifx\SetFigFont\undefined
% extract first six characters in \fmtname
\def\x#1#2#3#4#5#6#7\relax{\def\x{#1#2#3#4#5#6}}%
\expandafter\x\fmtname xxxxxx\relax \def\y{splain}%
\ifx\x\y   % LaTeX or SliTeX?
\gdef\SetFigFont#1#2#3{%
  \ifnum #1<17\tiny\else \ifnum #1<20\small\else
  \ifnum #1<24\normalsize\else \ifnum #1<29\large\else
  \ifnum #1<34\Large\else \ifnum #1<41\LARGE\else
     \huge\fi\fi\fi\fi\fi\fi
  \csname #3\endcsname}%
\else
\gdef\SetFigFont#1#2#3{\begingroup
  \count@#1\relax \ifnum 25<\count@\count@25\fi
  \def\x{\endgroup\@setsize\SetFigFont{#2pt}}%
  \expandafter\x
    \csname \romannumeral\the\count@ pt\expandafter\endcsname
    \csname @\romannumeral\the\count@ pt\endcsname
  \csname #3\endcsname}%
\fi
\fi\endgroup
\begin{picture}(4094,2594)(2229,-2483)
\put(3976,-736){\makebox(0,0)[lb]{\smash{\SetFigFont{12}{14.4}{rm}$e$}}}
\put(4126,-1786){\makebox(0,0)[lb]{\smash{\SetFigFont{12}{14.4}{rm}$P_W$}}}
\end{picture}

\]
\ni with $e$ being a matrix unit. The string containing $e$ enters
the box for $p_w$ either connecting two distinct $f$'s or two
strings of the same one. In the first case the result is zero since
$e$ belongs to precisely one of the direct summands.  In the second
case it is zero because of the properties of the $f$'s.

Condition (ii) follows similarly, noting that a picture like the
above will occur unless all the strings of $x$ are through strings.

For condition (iii) we observe that the ideal $I_k$ is linearly
spanned by $x$'s as above with less than $k$ through-strings, so we
may suppose $x$ is composed of through-strings, each with a matrix
unit in its box. Note that any relation true in $B_p$ is true in
$P_k$. Hence if $j\geq\l(p)$, $f_j=0$ and the identity of $B_p$ is
a linear combination of tangles with less than $\l(p)$
through-strings. So we can suppose that in $x$ there is no sequence
of $\l(p)$ strings in a row whose matrix unit labels are in the
same simple  summand as $p$. Thus by multiplying $x$ to the left
and right by tangles with the appropriate matrix unit labels, we
get $axb=p_w$ for some word $w$ in $\frak S(\cal M)$ of length $k$.
Thus (iii) will follow provided (iv) holds.

We calculate the normalized trace of $p_w$. It is 
\[
	\begin{picture}(0,0)%
\epsfig{file=ps/pic73.ps}%
\end{picture}%
\setlength{\unitlength}{0.00041700in}%
%
\begingroup\makeatletter\ifx\SetFigFont\undefined
% extract first six characters in \fmtname
\def\x#1#2#3#4#5#6#7\relax{\def\x{#1#2#3#4#5#6}}%
\expandafter\x\fmtname xxxxxx\relax \def\y{splain}%
\ifx\x\y   % LaTeX or SliTeX?
\gdef\SetFigFont#1#2#3{%
  \ifnum #1<17\tiny\else \ifnum #1<20\small\else
  \ifnum #1<24\normalsize\else \ifnum #1<29\large\else
  \ifnum #1<34\Large\else \ifnum #1<41\LARGE\else
     \huge\fi\fi\fi\fi\fi\fi
  \csname #3\endcsname}%
\else
\gdef\SetFigFont#1#2#3{\begingroup
  \count@#1\relax \ifnum 25<\count@\count@25\fi
  \def\x{\endgroup\@setsize\SetFigFont{#2pt}}%
  \expandafter\x
    \csname \romannumeral\the\count@ pt\expandafter\endcsname
    \csname @\romannumeral\the\count@ pt\endcsname
  \csname #3\endcsname}%
\fi
\fi\endgroup
\begin{picture}(10200,2129)(826,-2176)
\put(5176,-1636){\makebox(0,0)[lb]{\smash{\SetFigFont{12}{14.4}{rm}$\cdots$}}}
\put(7426,-1636){\makebox(0,0)[lb]{\smash{\SetFigFont{12}{14.4}{rm}$)$}}}
\put(7726,-1636){\makebox(0,0)[lb]{\smash{\SetFigFont{12}{14.4}{rm}$= \delta^{-k}\Pi_iZ($}}}
\put(11026,-1636){\makebox(0,0)[lb]{\smash{\SetFigFont{12}{14.4}{rm}$)$}}}
\put(826,-1636){\makebox(0,0)[lb]{\smash{\SetFigFont{12}{14.4}{rm}$\delta^{-k}Z($}}}
\put(2401,-736){\makebox(0,0)[lb]{\smash{\SetFigFont{12}{14.4}{rm}$p_1$}}}
\put(3901,-736){\makebox(0,0)[lb]{\smash{\SetFigFont{12}{14.4}{rm}$p_2$}}}
\put(6301,-736){\makebox(0,0)[lb]{\smash{\SetFigFont{12}{14.4}{rm}$p_r$}}}
\put(6151,-1636){\makebox(0,0)[lb]{\smash{\SetFigFont{12}{14.4}{rm}$f_{m_r}$}}}
\put(3826,-1636){\makebox(0,0)[lb]{\smash{\SetFigFont{12}{14.4}{rm}$f_{m_2}$}}}
\put(2251,-1636){\makebox(0,0)[lb]{\smash{\SetFigFont{12}{14.4}{rm}$f_{m_1}$}}}
\put(9976,-736){\makebox(0,0)[lb]{\smash{\SetFigFont{12}{14.4}{rm}$p_i$}}}
\put(9826,-1636){\makebox(0,0)[lb]{\smash{\SetFigFont{12}{14.4}{rm}$f_{m_i}$}}}
\end{picture}

\]
Now the partition function on $P_m$, restricted to $B_p$,
gives a Markov trace which will be normalized after division by
$Z(\blright{p} )^m=(\d\tau_p)^m$. So
$Z(\blright{f_m})=\d^m\tau_p^m T_m(\tau_{p}\d)$.
Hence tr$(p_w)=\prod^r_{i=1} \tau_{p_i}T_{m_i}(\tau_{p_i}\d)$.

Finally we must calculate the multiplicities dim$(p_vP_kp_w)$ for
$v$ of length $k$ and $w$ of length $k-1$. We must consider
diagrams of the form
\[
	\begin{picture}(0,0)%
\epsfig{file=ps/pic74.ps}%
\end{picture}%
\setlength{\unitlength}{0.00041700in}%
%
\begingroup\makeatletter\ifx\SetFigFont\undefined
% extract first six characters in \fmtname
\def\x#1#2#3#4#5#6#7\relax{\def\x{#1#2#3#4#5#6}}%
\expandafter\x\fmtname xxxxxx\relax \def\y{splain}%
\ifx\x\y   % LaTeX or SliTeX?
\gdef\SetFigFont#1#2#3{%
  \ifnum #1<17\tiny\else \ifnum #1<20\small\else
  \ifnum #1<24\normalsize\else \ifnum #1<29\large\else
  \ifnum #1<34\Large\else \ifnum #1<41\LARGE\else
     \huge\fi\fi\fi\fi\fi\fi
  \csname #3\endcsname}%
\else
\gdef\SetFigFont#1#2#3{\begingroup
  \count@#1\relax \ifnum 25<\count@\count@25\fi
  \def\x{\endgroup\@setsize\SetFigFont{#2pt}}%
  \expandafter\x
    \csname \romannumeral\the\count@ pt\expandafter\endcsname
    \csname @\romannumeral\the\count@ pt\endcsname
  \csname #3\endcsname}%
\fi
\fi\endgroup
\begin{picture}(1644,3690)(1479,-3274)
\put(1951,-436){\makebox(0,0)[lb]{\smash{\SetFigFont{12}{14.4}{rm}$P_W$}}}
\put(2101,-1486){\makebox(0,0)[lb]{\smash{\SetFigFont{12}{14.4}{rm}$x$}}}
\put(2026,-2686){\makebox(0,0)[lb]{\smash{\SetFigFont{12}{14.4}{rm}$P_V$}}}
\put(2101,-961){\makebox(0,0)[lb]{\smash{\SetFigFont{12}{14.4}{rm}$\cdots$}}}
\put(2101,-2161){\makebox(0,0)[lb]{\smash{\SetFigFont{12}{14.4}{rm}$\cdots$}}}
\put(2101,-3211){\makebox(0,0)[lb]{\smash{\SetFigFont{12}{14.4}{rm}$\cdots$}}}
\put(2101,164){\makebox(0,0)[lb]{\smash{\SetFigFont{12}{14.4}{rm}$\cdots$}}}
\end{picture}

\]
\ni where $x$ is a Temperley Lieb diagram decorated with matrix
units as before. Arguing on $p_v$, \ $x$ has only through-strings,
$v=w_p$ for some $p\in{\cal M}$, and the first $k-1$ strings of $x$
are labelled by the elements of $\cal M$ in $w$. Thus the diagram
is in fact equal to
\[
	\begin{picture}(0,0)%
\epsfig{file=ps/pic75.ps}%
\end{picture}%
\setlength{\unitlength}{0.00041700in}%
%
\begingroup\makeatletter\ifx\SetFigFont\undefined
% extract first six characters in \fmtname
\def\x#1#2#3#4#5#6#7\relax{\def\x{#1#2#3#4#5#6}}%
\expandafter\x\fmtname xxxxxx\relax \def\y{splain}%
\ifx\x\y   % LaTeX or SliTeX?
\gdef\SetFigFont#1#2#3{%
  \ifnum #1<17\tiny\else \ifnum #1<20\small\else
  \ifnum #1<24\normalsize\else \ifnum #1<29\large\else
  \ifnum #1<34\Large\else \ifnum #1<41\LARGE\else
     \huge\fi\fi\fi\fi\fi\fi
  \csname #3\endcsname}%
\else
\gdef\SetFigFont#1#2#3{\begingroup
  \count@#1\relax \ifnum 25<\count@\count@25\fi
  \def\x{\endgroup\@setsize\SetFigFont{#2pt}}%
  \expandafter\x
    \csname \romannumeral\the\count@ pt\expandafter\endcsname
    \csname @\romannumeral\the\count@ pt\endcsname
  \csname #3\endcsname}%
\fi
\fi\endgroup
\begin{picture}(1844,2410)(1479,-3349)
\put(2026,-2686){\makebox(0,0)[lb]{\smash{\SetFigFont{12}{14.4}{rm}$P_V$}}}
\put(2101,-1936){\makebox(0,0)[lb]{\smash{\SetFigFont{12}{14.4}{rm}$\cdots$}}}
\put(2101,-3286){\makebox(0,0)[lb]{\smash{\SetFigFont{12}{14.4}{rm}$\cdots$}}}
\put(2926,-1636){\makebox(0,0)[lb]{\smash{\SetFigFont{12}{14.4}{rm}$e$}}}
\end{picture}

\]
\ni where $e$ is a matrix unit with one subscript fixed.
These diagrams span $p_v P_k p_w$ and the sesquilinear form given
by $(x,y)={\op{tr}}(y^*x)$ is diagonal with non-zero entries.
Hence dim$(p_vP_kp_w)=n_p$. \qed

\bs\ni{\bf 3.2. Duality}

If $P=\cup_k P_k$ is  a planar algebra, the filtered algebras
$\l_n(P)$, where $\l_n(P)_k=P_{n,n+k}$, for fixed $n$, have
natural planar algebra structures. For $n$ even, this is rather
obvious --- just add $n$ straight vertical lines to the left of a
tangle. But if $n$ is odd one must be more careful because of
orientations. In fact $\l_1(P)$ and $P$ are not isomorphic in
general, even as filtered algebras, as one can see from example
2.9. We begin by describing the planar algebra structure on
$\l_1(P)$.

If $T$ is an unlabelled $k$-tangle we define the unlabelled
$(k+1)$-tangle $\tilde T$ to be the tangle consisting of a vertical
straight line from (1,0) to (1,1), and the tangle $T$, with all its
orientations reversed, shifted by 1 in the positive $x$ direction.
Also in $\tilde T$ each internal $p$-box is replaced by a
$(p+1)$-box with the first and last distinguished boundary points
connected by a short curve.  The procedure is illustrated in Figure 3.2.1.
%%%%%%%%%%%%%%%%%%%%%%%%%%%%%%%%%%%%%%%%%%%%%%%%%%%%%%%%
\[
	\begin{picture}(0,0)%
\epsfig{file=ps/pic76.ps}%
\end{picture}%
\setlength{\unitlength}{0.00041700in}%
%
\begingroup\makeatletter\ifx\SetFigFont\undefined
% extract first six characters in \fmtname
\def\x#1#2#3#4#5#6#7\relax{\def\x{#1#2#3#4#5#6}}%
\expandafter\x\fmtname xxxxxx\relax \def\y{splain}%
\ifx\x\y   % LaTeX or SliTeX?
\gdef\SetFigFont#1#2#3{%
  \ifnum #1<17\tiny\else \ifnum #1<20\small\else
  \ifnum #1<24\normalsize\else \ifnum #1<29\large\else
  \ifnum #1<34\Large\else \ifnum #1<41\LARGE\else
     \huge\fi\fi\fi\fi\fi\fi
  \csname #3\endcsname}%
\else
\gdef\SetFigFont#1#2#3{\begingroup
  \count@#1\relax \ifnum 25<\count@\count@25\fi
  \def\x{\endgroup\@setsize\SetFigFont{#2pt}}%
  \expandafter\x
    \csname \romannumeral\the\count@ pt\expandafter\endcsname
    \csname @\romannumeral\the\count@ pt\endcsname
  \csname #3\endcsname}%
\fi
\fi\endgroup
\begin{picture}(10828,4309)(1179,-4648)
\put(2392,-4576){\makebox(0,0)[lb]{\smash{\SetFigFont{12}{14.4}{rm}The tangle $T$.}}}
\put(8377,-4561){\makebox(0,0)[lb]{\smash{\SetFigFont{12}{14.4}{rm}The tangle ${\tilde T}$.}}}
\put(2790,-2210){\makebox(0,0)[lb]{\smash{\special{ps:gsave currentpoint currentpoint translate
-315.0 rotate neg exch neg exch translate}\SetFigFont{12}{14.4}{rm}$*$\special{ps:currentpoint grestore moveto}}}}
\put(4105,-2939){\makebox(0,0)[lb]{\smash{\special{ps:gsave currentpoint currentpoint translate
-15.0 rotate neg exch neg exch translate}\SetFigFont{12}{14.4}{rm}$*$\special{ps:currentpoint grestore moveto}}}}
\put(4790,-1555){\makebox(0,0)[lb]{\smash{\special{ps:gsave currentpoint currentpoint translate
-30.0 rotate neg exch neg exch translate}\SetFigFont{12}{14.4}{rm}$*$\special{ps:currentpoint grestore moveto}}}}
\put(10119,-2985){\makebox(0,0)[lb]{\smash{\special{ps:gsave currentpoint currentpoint translate
-15.0 rotate neg exch neg exch translate}\SetFigFont{12}{14.4}{rm}$*$\special{ps:currentpoint grestore moveto}}}}
\put(8819,-2286){\makebox(0,0)[lb]{\smash{\special{ps:gsave currentpoint currentpoint translate
-315.0 rotate neg exch neg exch translate}\SetFigFont{12}{14.4}{rm}$*$\special{ps:currentpoint grestore moveto}}}}
\put(10801,-1561){\makebox(0,0)[lb]{\smash{\special{ps:gsave currentpoint currentpoint translate
-30.0 rotate neg exch neg exch translate}\SetFigFont{12}{14.4}{rm}$*$\special{ps:currentpoint grestore moveto}}}}
\end{picture}

\]
\begin{center}
Figure 3.2.1
\end{center}
\bs
To each internal box $B$ of $T$ there corresponds in the obvious
way a box $\tilde B$ of $\tilde T$. If $T$ is labelled by
$L=\coprod_{k > 0} L_k$, \
$\tilde T$ will be given the obvious labelling by $\tilde L$, \
$\tilde L_k= L_{k-1}(\tilde L_1=\emptyset)$.

\bs\ni
{\bf Proposition 3.2.2} Let $P=\cup_k P_k$ be a planar algebra
with parameters $\d_1=Z(\caright)$ and
$\d_2=Z(\caleft)$. Assume $P$ is presented on itself by
$\Phi$. Then $\l_1(P)$ is a planar algebra with parameters
$\d_2,\d_1$, presented on itself by $\l_1(\Phi)$ where $\l_1(\Phi)$
is defined by $\l_1(\Phi)(T)=\d^{-p}\Phi(\tilde T)$, \ $p$ being
the number of internal boxes in $T$.

\bs
{\sc Proof.} First note how  the labels in a tangle of
${\cal P}(\l_1(P))$ give valid labels for ${\cal P}(P)$ because of the
inclusion $\l_1(P)_k\subset P_{k+1}$. That $\l_1(\Phi)$ is a
filtered algebra homomorphism is obvious. The annular invariance of
ker $\l_1(\Phi)$ follows immediately from that of $\Phi$, by
representing elements of $P_{1,1+k}$ as linear combinations of
tangles with vertical first string and applying $\sim$ to linear
combinations. Thus $\l_1(P)$ is a general planar algebra.

Now $\l_1(P)_0=P_{1,1}$ and $\l_1(P)_{1,1}=P_{2,2}$ which has
dimension 1 since $P$ is a planar algebra. So
$\l_1(P)$ is planar. The multiplicativity property for $\l_1(P)$
follows immediately from that of $P$, where orientations on
networks without boxes are reversed. \qed

\bs In the next two lemmas, $A$ will be the ${\cal A}(\emptyset)$
element
\[
	\begin{picture}(0,0)%
\epsfig{file=ps/pic77.ps}%
\end{picture}%
\setlength{\unitlength}{0.00041700in}%
%
\begingroup\makeatletter\ifx\SetFigFont\undefined
% extract first six characters in \fmtname
\def\x#1#2#3#4#5#6#7\relax{\def\x{#1#2#3#4#5#6}}%
\expandafter\x\fmtname xxxxxx\relax \def\y{splain}%
\ifx\x\y   % LaTeX or SliTeX?
\gdef\SetFigFont#1#2#3{%
  \ifnum #1<17\tiny\else \ifnum #1<20\small\else
  \ifnum #1<24\normalsize\else \ifnum #1<29\large\else
  \ifnum #1<34\Large\else \ifnum #1<41\LARGE\else
     \huge\fi\fi\fi\fi\fi\fi
  \csname #3\endcsname}%
\else
\gdef\SetFigFont#1#2#3{\begingroup
  \count@#1\relax \ifnum 25<\count@\count@25\fi
  \def\x{\endgroup\@setsize\SetFigFont{#2pt}}%
  \expandafter\x
    \csname \romannumeral\the\count@ pt\expandafter\endcsname
    \csname @\romannumeral\the\count@ pt\endcsname
  \csname #3\endcsname}%
\fi
\fi\endgroup
\begin{picture}(3644,2144)(1779,-2483)
\end{picture}

\]
\ni where the actual number of boundary points is as required by
context.

\bs\ni
{\bf Lemma 3.2.3} If $P$ is a planar algebra, $\pi_A$ defines
a linear isomorphism between $P_k$ and $\l_1(P)_k$, for each $k >0$.

\bigskip
{\sc Proof.} By a little isotopy and the definition of
$P_{1,k+1}$, \ $\pi_A$ is onto. But
\[
	\begin{picture}(0,0)%
\epsfig{file=ps/pic78.ps}%
\end{picture}%
\setlength{\unitlength}{0.00041700in}%
%
\begingroup\makeatletter\ifx\SetFigFont\undefined
% extract first six characters in \fmtname
\def\x#1#2#3#4#5#6#7\relax{\def\x{#1#2#3#4#5#6}}%
\expandafter\x\fmtname xxxxxx\relax \def\y{splain}%
\ifx\x\y   % LaTeX or SliTeX?
\gdef\SetFigFont#1#2#3{%
  \ifnum #1<17\tiny\else \ifnum #1<20\small\else
  \ifnum #1<24\normalsize\else \ifnum #1<29\large\else
  \ifnum #1<34\Large\else \ifnum #1<41\LARGE\else
     \huge\fi\fi\fi\fi\fi\fi
  \csname #3\endcsname}%
\else
\gdef\SetFigFont#1#2#3{\begingroup
  \count@#1\relax \ifnum 25<\count@\count@25\fi
  \def\x{\endgroup\@setsize\SetFigFont{#2pt}}%
  \expandafter\x
    \csname \romannumeral\the\count@ pt\expandafter\endcsname
    \csname @\romannumeral\the\count@ pt\endcsname
  \csname #3\endcsname}%
\fi
\fi\endgroup
\begin{picture}(3644,2149)(1779,-2183)
\end{picture}

\]
\ni provides an inverse for $\pi_A$, up to a non-zero scalar.
So $\pi_A$ is an isomorphism. \qed


{\bf Lemma 3.2.4} The subset $S$ of the planar algebra
generates $P$ as a planar algebra iff $\pi_A(S)$ generates
$\l_1(P)$ as a planar algebra.

\bigskip
{\sc Proof.} $(\Rightarrow)$ \  Given a tangle $T$ in
${\cal P}_{1,k+1}(S)$, it suffices to exhibit a tangle $T_A$
in ${\cal P}_k(\pi_A(S))$ with $\Phi(\tilde T_A)$ being a multiple of
$\Phi(T)$. We create $T_A$ from $T$ by eliminating the first
string, reversing all orientations and otherwise changing only in
small neighborhoods of the internal boxes of $G$, sending a box
labelled $R\in S$ in $T$ to the box labelled $\pi_A(R)$ in $T_A$ as
below:
\[
	\begin{picture}(0,0)%
\epsfig{file=ps/pic79.ps}%
\end{picture}%
\setlength{\unitlength}{0.00041700in}%
%
\begingroup\makeatletter\ifx\SetFigFont\undefined
% extract first six characters in \fmtname
\def\x#1#2#3#4#5#6#7\relax{\def\x{#1#2#3#4#5#6}}%
\expandafter\x\fmtname xxxxxx\relax \def\y{splain}%
\ifx\x\y   % LaTeX or SliTeX?
\gdef\SetFigFont#1#2#3{%
  \ifnum #1<17\tiny\else \ifnum #1<20\small\else
  \ifnum #1<24\normalsize\else \ifnum #1<29\large\else
  \ifnum #1<34\Large\else \ifnum #1<41\LARGE\else
     \huge\fi\fi\fi\fi\fi\fi
  \csname #3\endcsname}%
\else
\gdef\SetFigFont#1#2#3{\begingroup
  \count@#1\relax \ifnum 25<\count@\count@25\fi
  \def\x{\endgroup\@setsize\SetFigFont{#2pt}}%
  \expandafter\x
    \csname \romannumeral\the\count@ pt\expandafter\endcsname
    \csname @\romannumeral\the\count@ pt\endcsname
  \csname #3\endcsname}%
\fi
\fi\endgroup
\begin{picture}(4987,1980)(1179,-1801)
\put(1651,-1636){\makebox(0,0)[lb]{\smash{\SetFigFont{12}{14.4}{rm}$\cdots$}}}
\put(1651,-211){\makebox(0,0)[lb]{\smash{\SetFigFont{12}{14.4}{rm}$\cdots$}}}
\put(4501,-811){\makebox(0,0)[lb]{\smash{\SetFigFont{12}{14.4}{rm}$\pi_A(R)$}}}
\put(1651,-811){\makebox(0,0)[lb]{\smash{\SetFigFont{12}{14.4}{rm}$R$}}}
\put(4696,-211){\makebox(0,0)[lb]{\smash{\SetFigFont{12}{14.4}{rm}$\cdots$}}}
\put(4681,-1621){\makebox(0,0)[lb]{\smash{\SetFigFont{12}{14.4}{rm}$\cdots$}}}
\end{picture}

\]
Then by definition $\tilde T_A$ will be exactly like $T$ except
near its boxes where it will look as below:
\[
	\begin{picture}(0,0)%
\epsfig{file=ps/pic80.ps}%
\end{picture}%
\setlength{\unitlength}{0.00041700in}%
%
\begingroup\makeatletter\ifx\SetFigFont\undefined
% extract first six characters in \fmtname
\def\x#1#2#3#4#5#6#7\relax{\def\x{#1#2#3#4#5#6}}%
\expandafter\x\fmtname xxxxxx\relax \def\y{splain}%
\ifx\x\y   % LaTeX or SliTeX?
\gdef\SetFigFont#1#2#3{%
  \ifnum #1<17\tiny\else \ifnum #1<20\small\else
  \ifnum #1<24\normalsize\else \ifnum #1<29\large\else
  \ifnum #1<34\Large\else \ifnum #1<41\LARGE\else
     \huge\fi\fi\fi\fi\fi\fi
  \csname #3\endcsname}%
\else
\gdef\SetFigFont#1#2#3{\begingroup
  \count@#1\relax \ifnum 25<\count@\count@25\fi
  \def\x{\endgroup\@setsize\SetFigFont{#2pt}}%
  \expandafter\x
    \csname \romannumeral\the\count@ pt\expandafter\endcsname
    \csname @\romannumeral\the\count@ pt\endcsname
  \csname #3\endcsname}%
\fi
\fi\endgroup
\begin{picture}(7912,1980)(652,-1802)
\put(1877,-812){\makebox(0,0)[lb]{\smash{\SetFigFont{12}{14.4}{rm}$\pi_A(R)$}}}
\put(2072,-212){\makebox(0,0)[lb]{\smash{\SetFigFont{12}{14.4}{rm}$\cdots$}}}
\put(2057,-1622){\makebox(0,0)[lb]{\smash{\SetFigFont{12}{14.4}{rm}$\cdots$}}}
\put(6874,-214){\makebox(0,0)[lb]{\smash{\SetFigFont{12}{14.4}{rm}$\cdots$}}}
\put(6859,-1624){\makebox(0,0)[lb]{\smash{\SetFigFont{12}{14.4}{rm}$\cdots$}}}
\put(6902,-812){\makebox(0,0)[lb]{\smash{\SetFigFont{12}{14.4}{rm}$R$}}}
\put(4173,-586){\makebox(0,0)[lb]{\smash{\SetFigFont{12}{14.4}{rm}$\Phi$}}}
\end{picture}

\]
\ni Thus $\Phi(\tilde T_A)$ is a multiple of $\Phi(T)$.

$(\Leftarrow)$ \ Given $x\in P$, \ $\pi_A(\tangle{x})$
is in $\l_1(P)$ so by hypothesis it is the image under
$\l_1(\Phi)$ of a linear combination of tangles  labelled by
elements of $\pi_A(S)$ which are in turn images under $\Phi$ of
tangles  labelled by
elements of $S$ (up to nonzero scalars). Using the tangle of Lemma
3.2.3 to invert $\pi_A$, we are done. \qed

\bs By iterating the procedure $P\mapsto\l_1(P)$, we see that all
the $\l_n(P)$ have natural planar algebra structures, but observe
that all the $\l_{2n}(P)$ are isomorphic to $P$ as planar algebras
via the endomorphism (often called ``le shift de deux") defined by
adding two straight vertical strings to the left of a tangle. We
leave the details to the reader.

The planar algebra $\l_1(P)$ is said to be the {\it dual}
of the planar algebra $P$, and we have $\l_1(\l_1(P))\simeq P$
as planar algebras.

In the case of Example 2.9, $P^G_2$ is the group algebra
$\Bbb C G$ and $\l_1(P^G)_2$ is $\ell^\i(G)$. The tangle $\pi_A$
gives a {\it linear} isomorphism between the two. Thus planar
algebra duality extends the duality between a finite group and its
dual object.




